\section{Sharp boundaries and impossibility arguments}
\label{sec:barriers}

A general framework is trustworthy only if it explains what it cannot encode.
The following barriers show that several recurring difficulties are
mathematical, not deficiencies of notation.

\subsection{No total universal primitive-recursive evaluator}

\begin{theorem}[Diagonal obstruction]
\label{thm:no-universal-pr}
There is no binary primitive-recursive function \(U(e,x)\) such that every
unary primitive-recursive function \(f\) has an index \(e\) satisfying
\(f(x)=U(e,x)\) for all \(x\).
\end{theorem}

\begin{proof}
If such \(U\) existed, then
\[
 g(n)=U(n,n)+1
\]
would be primitive recursive.  By universality, choose \(e\) with
\(g(n)=U(e,n)\) for every \(n\).  At \(n=e\),
\[
 U(e,e)=g(e)=U(e,e)+1,
\]
a contradiction.
\end{proof}

Consequently, a first-order theorem cannot receive an arbitrary PR function
code and evaluate it by one total PR symbol.  Fixed evaluator schemas and
finite trace verifiers are not a temporary workaround; they are the maximal
uniform interface compatible with primitive recursiveness.

\subsection{A cardinality bound does not yield a height bound}

\begin{theorem}[No uniform conversion from ``at most one'' to a cutoff]
\label{thm:cardinality-not-height}
There is a primitive-recursive decidable predicate \(P(e,t)\) such that for
every \(e\) at most one \(t\) satisfies \(P(e,t)\), but there is no
primitive-recursive function \(H(e)\) satisfying
\[
 P(e,t)\to t\le H(e)
\]
for all \(e,t\).
\end{theorem}

\begin{proof}
Let \(P(e,t)\) say that Turing machine \(e\), on blank input, halts for the
first time exactly at step \(t\).  Bounded deterministic simulation is
primitive recursive, so \(P\) is primitive-recursive and decidable.  A machine
has at most one first halting time.

If a primitive-recursive \(H\) bounded every possible first halting time, the
halting problem would be decidable: simulate machine \(e\) for \(H(e)\) steps.
If it has not halted by then, the implication says it never halts.  This
contradicts undecidability of the halting problem.
\end{proof}

Thus even a uniform cardinality bound of one carries no search-stopping
information.  The obstacle in the Faltings and Siegel packets is exactly the
missing location/height data, not the finite filtering algorithm.

\subsection{Exact sign selection and LLPO}

One standard form of LLPO says: for every binary sequence with at most one
entry equal to one, either all even entries are zero or all odd entries are
zero.  From such a sequence \(\alpha\), construct the regular real name
corresponding to
\[
 x=\sum_{n\ge0}(\alpha(2n)-\alpha(2n+1))2^{-n-1}.
\]
The at-most-one premise gives an explicit tail modulus.  If a uniform
procedure always decides \(x\ge0\) or \(x\le0\), then:
\begin{itemize}
\item \(x\ge0\) rules out an odd-position one;
\item \(x\le0\) rules out an even-position one.
\end{itemize}
When the sequence is all zero, both conclusions are valid.  Hence real sign
selection realizes LLPO.

The exact-axis and sharp-endpoint obstruction arguments in the inspected
packets reduce their selectors to this sign principle.  Replacing the real
encoding or choosing a different normal form cannot remove the obstruction.
Overlapping branches avoid it because they do not decide which side of an
exact boundary contains zero.

\subsection{Negative inequality is not positive apartness}

Intuitionistically, \(\neg(x=0)\) need not provide a rational \(d>0\) with
\(\abs{x}\ge d\).  Algorithms for inversion, normalization, or density need
the latter.  A principle converting every negative nonzeroness proof into a
positive modulus is a Markov-style principle for represented numbers.  The
framework therefore treats apartness witnesses as data and never derives them
from bare negation unless the input class has a specific finite argument.

For rationals and algebraic numbers with an isolating certificate, positive
apartness can be computed.  For arbitrary regular names, it cannot.

\subsection{Ineffective existence cannot be repaired by a verifier}

A finite verifier may prove that any supplied constant table implies an
asymptotic estimate.  This does not imply a primitive-recursive function
producing such a table.  The difference mirrors proof checking versus theorem
search: local verification may be elementary even when uniform production is
open or noncomputable.

Quantifier order detects the issue.  The useful statement is
\[
 \forall D\,\forall A\;\exists N\;\forall x\ge N\;E(D,A,x),
\]
with \(N\) produced before \(x\).  A family of proofs of
\(\forall D,A,x\;\exists N_x\le x\;E(D,A,x)\) is not a substitute.  The
finite-trace framework preserves the former order by making the threshold an
output of a certificate producer.

\subsection{Semantic truth is not automatic iPRA provability}

A primitive-recursive function can have true arithmetic properties whose
proofs require induction or consistency strength beyond a particular weak
theory.  Therefore the slogan ``the algorithm is PR, so iPRA proves it'' is
false.  The framework always asks for the local invariant derivations
\eqref{eq:T0}--\eqref{eq:T2}.  The finite-trace theorem composes proofs; it
does not create missing mathematical invariants.

\subsection{Maximal positive statement principle}

When an unrestricted theorem meets one of these barriers, the goal is not to
stop at ``impossible.''  One should determine the maximal positive family:
\begin{itemize}
\item approximate witness at every precision;
\item exact invariant or projective data;
\item exact witness under a positive gap;
\item certificate-relative theorem;
\item theorem for rational or finitely presented inputs;
\item obstruction showing why the omitted uniform selector is stronger.
\end{itemize}
The Euler--Rodrigues and Dirichlet--Kronecker packets already exemplify this
pattern.  The \(\CFP\) framework makes it a standard deliverable for every
future packet.
